%% kapa: exam=1 round=23 year=2012 period=2 form=A subject=appraisal_law subject_ko=부동산관계법규 seq=40 number=80
%% source: Q-Net 감정평가사 기출문제 (cst00309) artlSeq=5052580 file=raw/1008/5052580/2012년 제23회 감정평가사 제1차시험 기출문제(2교시 A형).pdf
%% conversion: kordoc 4.24.51 -> tools/gen1.py (auto); review=proofread-sonnet; flags=none
\begin{question}{80}
부동산등기법상 신탁등기에 관한 설명으로 옳지 않은 것은?
\begin{choices}
  \item 수익자 또는 위탁자는 수탁자를 대위하여 신탁등기를 신청할 수 있다.
  \item 수탁자가 여러 명인 경우 등기관은 신탁재산이 공유인 뜻을 기록하여야 한다.
  \item 신탁등기의 신청은 해당 신탁으로 인한 권리의 이전 또는 보존이나 설정등기의 신청과 동시에 하여야 한다.
  \item 신탁을 원인으로 위탁자가 자기 명의의 재산을 수탁자에게 처분하는 경우 신탁등기의 신청은 위탁자를 등기의무자로 하고 수탁자를 등기권리자로 한다.
  \item 신탁재산에 속한 권리가 소멸된 경우 신탁등기의 말소신청은 신탁된 권리의 말소등기의 신청과 동시에 하여야 한다.
\end{choices}
\end{question}
\answer{2}
